% ============================================================================
%  fig_pipeline_overview.tex -- whole-project dataflow, P0 through P5.
%  Requires preamble_figures.tex in the document preamble.
%
%  REWORK (2026-08-30, owner directive): the first draft of this figure gave
%  P2's empty whitelist and its "never read" dead-edge the dominant visual
%  weight (a failnode, a separate artifact node, AND a diagonal red edge
%  snaking across the whole canvas) -- disproportionate to one diagnostic
%  result out of five phases of findings. P2 is real and stays, but as ONE
%  compact node like every other phase; the whitelist artifact and its
%  dead-edge move entirely into fig_p2_reliability.tex, the dedicated P2
%  figure, where that level of detail belongs. The findings banner is
%  replaced: instead of restating P2+P3's shared negative result, it now
%  carries the project's actual throughline, the dissociation that P2, P3,
%  P4 and P5 jointly establish -- fractal features are redundant
%  in-distribution, ACTIVELY HURT under 2 of 6 real attacks, and are the
%  only source of a real, significant advantage under genuine cross-corpus
%  transfer. That is the finding the whole pipeline was built to surface,
%  not any single phase's result in isolation.
%
%  LAYOUT: verified by rendering, not by eye. Node gaps are >= 1.9cm so
%  edge labels sit in clear space; the findings banner spans the lower band
%  where nothing else is drawn.
% ============================================================================
\begin{tikzpicture}[x=1cm,y=1cm]

% ------------------------------------------------------------ input corpora
\node[corpus, minimum width=28mm] (gagle) at (2.4,8.3)
  {\textbf{GAGLE}\\[1pt] 223{,}834 docs\\ 5 domains $\times$ 8 prompt types};

\node[corpus, minimum width=28mm] (raid) at (17.7,1.9)
  {\textbf{RAID}\\[1pt] 68{,}212 docs\\ 11 generators, 8 domains};

% --------------------------------------------------------------- main spine
\node[phase, minimum width=30mm] (p0) at (2.4,6.0)
  {\textbf{P0}~~Corpus baseline\\[1pt]
   $\tau\!\in\![8,320]$, $\varepsilon\!=\!10^{-2}$\\
   drop 64 / clip 464 / min 320};

\node[phase, minimum width=32mm] (p1) at (7.6,6.0)
  {\textbf{P1}~~Feature store\\[1pt]
   Family A (54) $+$ B (17)\\ $+$ C (14) $=$ 85 features};

\node[phase, minimum width=32mm] (p3) at (12.6,6.0)
  {\textbf{P3}~~Binary detection\\[1pt]
   5 configs $\times$ 8 folds\\ $\times$ 12 arms $=$ 304 cells};

\node[evalnode, minimum width=28mm] (p4) at (17.7,7.9)
  {\textbf{P4}~~Robustness\\[1pt]
   6 attacks $\times$ 5 detectors\\ 39{,}810 edited docs};

\node[evalnode, minimum width=28mm] (p5) at (17.7,4.7)
  {\textbf{P5}~~Transfer $+$ multiclass\\[1pt]
   T2 / T3 / GAGLE$\rightarrow$RAID\\ 11-class generator ID};

% ------------------------------------------------------- P2: one node, not a spine
\node[phase, minimum width=34mm, text width=40mm] (p2) at (7.6,3.15)
  {\textbf{P2}~~Split-half reliability\\[1pt]
   pre-registered bar $\rho\!\geq\!0.60$: \textbf{0/20 pass}\\
   best feature $\rho = 0.063$};

% ------------------------------------------------------------------- edges
\draw[flow] (gagle) -- (p0);
\draw[flow] (p0) -- node[edgelab, pos=0.5, above=1pt]{single pass} (p1);
\draw[flow] (p1) -- node[edgelab, pos=0.5, above=1pt]{85 feat.} (p3);
% route the fan-out on a fixed vertical at x=15.2: it must clear P3's
% east edge (~14.2) AND P4/P5's west edges (~15.7). An earlier revision
% used ++(0.55,0), putting the connector INSIDE P4/P5 so the arrowheads
% clipped their first characters -- caught by rendering at 150 dpi.
\draw[flow] (p3.east) -- (15.2,6.0) |- (p4.west);
\draw[flow] (p3.east) -- (15.2,6.0) |- (p5.west);
\draw[flow] (raid) -- (p5);

\draw[flow] (p1.south) -- (p2.north);
\draw[flow] (p2.south) -- ++(0,-0.75)
  node[edgelab, pos=1, below, text width=44mm, align=center]
  {feeds P3's fractal family regardless of gate outcome ---\\
   see fig\_p2\_reliability.tex for the whitelist detail};

% --------------------------------------------- the project's dissociation (finding)
\node[finding, minimum width=112mm, text width=128mm] (spine) at (9.5,-1.55)
  {\textbf{The dissociation, established across four phases.}\quad
   \textbf{P2:} per-document fractal estimates barely agree with themselves
   ($\rho = 0.063$) --- domain-level noise, not document-level signal.\quad
   \textbf{P3:} consistent with that, dropping all 54 fractal features costs
   \emph{nothing} in-distribution (AUROC $0.9984$ either way) and wins on
   7/8 cross-domain/cross-generator folds.\quad
   \textbf{P4:} under real attacks, dropping fractal features
   \textcolor{okGreen!60!black}{\textbf{helps}} on the two that matter most
   (sentence-reorder, typo-noise) --- the opposite of a robustness benefit.\quad
   \textbf{P5:} the one place fractal features show a real,
   population-matched advantage is genuine cross-\emph{corpus} transfer to
   RAID ($\Delta$AUROC $= 0.024$, 95\% CI excludes 0) --- a different kind
   of shift than P3's in-corpus folds or P4's edits.};

% ------------------------------------------------------------- group labels
% x=-1.7, not -0.9: at -0.9 "Construction" (widest label) ran right up to
% and under P0's second line ("tau in [8,320]..."), which grows past its
% 30mm minimum width -- caught by rendering.
\node[grouplab, anchor=west] at (-1.7,8.3) {Corpora};
\node[grouplab, anchor=west] at (-1.7,6.0) {Construction};
\node[grouplab, anchor=west] at (-1.7,3.15) {Validation};
\node[grouplab, anchor=west] at (-1.7,-1.55) {Findings};

\end{tikzpicture}
